% !TeX program = pdflatex
% =============================================================================
% Verlaufsplanung: Gewichtskraft & Federkraft (Doppellektion) -- Lektionsablauf-Tabelle
% UZH Praktikum I -- Sara Romer Urban, Kantonsschule im Lee, HS2026
% Klasse: 2e (26.08.2026)
%
% Kurzfassung von gewichtskraft-federkraft-lektionsplan.tex: nur Zeitplan und
% die wichtigsten Eckdaten, ohne Begruendung/Offene-Punkte-Diskussion. Layout
% an Zeitplan_Vorlage.tex angelehnt (longtable, alternierende Zeilenfarben),
% Spaltenwahl an kondensatoren-verlaufsplanung.tex angelehnt: Inhalt, Zeit,
% Lehrperson (Taetigkeit), SuS (Taetigkeit), Material, Sozialform.
% =============================================================================
\documentclass[11pt,a4paper]{article}

\usepackage[landscape,margin=1.6cm]{geometry}
\usepackage[ngerman]{babel}
\usepackage[T1]{fontenc}
\usepackage[utf8]{inputenc}
\usepackage{lmodern}
\usepackage{array}
\usepackage{longtable}
\usepackage{booktabs}
\usepackage[table]{xcolor}
\usepackage{enumitem}
\usepackage{amsmath}
\usepackage{fancyhdr}
\usepackage{hyperref}

\pagestyle{fancy}
\fancyhf{}
\setlength{\headheight}{14pt}
\lhead{Physik -- Gewichtskraft \& Federkraft}
\rhead{Lektionsablauf}
\cfoot{\thepage}

\setlength{\parindent}{0pt}
\setlength{\parskip}{4pt}

\definecolor{headgray}{RGB}{230,230,230}
\definecolor{lightrow}{RGB}{245,245,245}

\hypersetup{
  colorlinks=true,
  urlcolor=blue,
  linkcolor=blue,
  pdftitle={Lektionsablauf: Gewichtskraft \& Federkraft},
  pdfauthor={Loris Keller}
}

\begin{document}

\begin{center}
    {\LARGE \textbf{Lektionsablauf: Gewichtskraft \& Federkraft}}\\[0.3cm]
    \large Mechanik -- Doppellektion (2$\times$45 min) -- 2e, HS2026
\end{center}

\vspace{0.3cm}

\renewcommand{\arraystretch}{1.3}
\setlength{\LTpre}{6pt}
\setlength{\LTpost}{6pt}

\begin{longtable}{@{}p{5.3cm}p{1.3cm}p{5.9cm}p{5.9cm}p{3.5cm}p{2.4cm}@{}}
\rowcolor{headgray}
\textbf{Inhalt} & \textbf{Zeit} & \textbf{Lehrperson} & \textbf{Schülerinnen und Schüler} & \textbf{Material} & \textbf{Sozialform} \\
\midrule
\endfirsthead

\rowcolor{headgray}
\textbf{Inhalt} & \textbf{Zeit} & \textbf{Lehrperson} & \textbf{Schülerinnen und Schüler} & \textbf{Material} & \textbf{Sozialform} \\
\midrule
\endhead

\bottomrule
\endfoot

\rowcolor{lightrow}
\textbf{Wiederholung} \newline
Kraftformel \& Werkstatt-Rückblick: $F\sim a$, Kraftmesser-Grundprinzip.
&
6'\newline{\footnotesize Ist:\dotfill}
&
Fragt ab, sammelt Antworten, gibt kurze Denkpause (B1 Vergewisserungsphase).
&
Denken kurz allein/zu zweit nach, äussern Antworten im Plenum.
&
Tafel/Beamer
&
Plenum
\\

\textbf{Übung: Fussballschuss} \newline
$F=ma$ am neuen Beispiel (LZ1), Kinematik-Dynamik-Verknüpfung.
&
10'\newline{\footnotesize Ist:\dotfill}
&
Stellt Aufgabe, geht herum, moderiert Austausch im Plenum (C3 Think-Pair-Share).
&
Überlegen allein, tauschen sich zu zweit aus, rechnen im Arbeitsblatt.
&
Arbeitsblatt (Bild Fussballschuss)
&
Einzel- $\to$ Partnerarbeit $\to$ Plenum
\\

\rowcolor{lightrow}
\textbf{Aktivierung 1: Steinstoss} \newline
Gedankenaufgabe Erde vs. Mond -- Trägheit (LZ2).
&
4'\newline{\footnotesize Ist:\dotfill}
&
Stellt Frage, moderiert kurze Plenumsdiskussion (B1 Vergewisserungsphase).
&
Denken kurz allein nach, begründen im Arbeitsblatt, äussern Vermutung.
&
Arbeitsblatt
&
Einzelarbeit $\to$ Plenum
\\

\textbf{Aktivierung 2: Golfball} \newline
Gedankenaufgabe Golfball auf dem Mond -- Gewichtskraft (LZ2/LZ3).
&
4'\newline{\footnotesize Ist:\dotfill}
&
Moderiert Partnerdiskussion, bestätigt historisches Ergebnis erst im Plenum (kein Video).
&
Diskutieren zu zweit, begründen im Arbeitsblatt.
&
Arbeitsblatt (Bild Golfball)
&
Partnerarbeit $\to$ Plenum
\\

\rowcolor{lightrow}
\textbf{Theorie 1: Gewichtskraft} \newline
$F_G=m\cdot g$, Ortsfaktor $g$ für Erde und Mond.
&
7'\newline{\footnotesize Ist:\dotfill}
&
Erklärt an der Tafel, knapp gehalten.
&
Hören zu, füllen Theorie im Arbeitsblatt.
&
Tafel/Beamer, Arbeitsblatt
&
Lehrervortrag
\\

\textbf{Theorie 2: Federkraft-Konzept} \newline
Kräftegleichgewicht, Vorspannung, $\Delta x$ vs. $x$ (LZ6). Ohne Hookesches
Gesetz.
&
7'\newline{\footnotesize Ist:\dotfill}
&
Erklärt Kalibrierungsprinzip, zeigt Federkraftmesser-Bild.
&
Hören zu, füllen Theorie im Arbeitsblatt.
&
Tafel/Beamer (Bild Federkraftmesser), Arbeitsblatt
&
Lehrervortrag
\\

\rowcolor{lightrow}
\textbf{Übergang} \newline
Wechsel zur Partnerarbeit: Gruppen bilden, Stationsmaterial verteilen.
&
4'\newline{\footnotesize Ist:\dotfill}
&
Teilt Klasse in Paare ein, verteilt Material.
&
Bilden Paare, holen Material, bauen Station auf.
&
Massstäbe, Federn, Massen, Laptops
&
Organisation
\\

\midrule
\multicolumn{6}{@{}c@{}}{\textit{--- Ende Lektion 1 (42') / Beginn Lektion 2 ---}} \\
\midrule

\textbf{Lernaufgabe} \newline
Federkraftmesser-Experiment: SuS leiten $F=D\cdot\Delta x$ und $D$ selbst
her (LZ4, LZ5).
&
38'\newline{\footnotesize Ist:\dotfill}
&
Geht herum, gibt Hilfestellung, gibt die Formel nicht vor (D1 Lernaufgabe).
&
Messen zu zweit an zwei Federn, tragen Werte in Excel ein, lesen $D$ aus
der Trendlinie ab, bearbeiten das Aufgabenblatt.
&
Stationsaufbau (Massstab, 2 Federn, Massen), Laptop, Excel-Vorlage,
Aufgabenblatt
&
Partnerarbeit
\\

\rowcolor{lightrow}
\textbf{Übung: Einheitenumrechnung} \newline
N/cm $\to$ N/m, anhand eigener Messresultate (LZ7).
&
5'\newline{\footnotesize Ist:\dotfill}
&
Gibt kurzen Auftrag, sammelt Ergebnis im Plenum.
&
Rechnen individuell im Arbeitsblatt.
&
Arbeitsblatt, Taschenrechner
&
Einzelarbeit
\\

\textbf{Ausblick} \newline
Hinweis auf Aufgabenblatt für Vor-/Nachbereitung.
&
5'\newline{\footnotesize Ist:\dotfill}
&
Fasst zusammen, verweist auf das Aufgabenblatt.
&
Hören zu, stellen letzte Fragen.
&
--
&
Plenum
\\

\rowcolor{lightrow}
\multicolumn{1}{@{}p{5.3cm}}{\textbf{Zusammenfassung}} & \multicolumn{1}{p{1.3cm}}{\textbf{90' / 90'}} & \multicolumn{4}{@{}p{18.9cm}@{}}{\textit{Lektion 1: 42' (Phasen 1--7, inkl. Übergang). Lektion 2: 48' (Phasen 8--10). Bei Zeitdruck kürzen: zuerst Ausblick, danach Einheitenumrechnung (kann im Aufgabenblatt nachgeholt werden, da dort ohnehin an eigenen Messdaten geübt wird). Die Lernaufgabe nie kürzen, sie ist das didaktische Zentrum.}} \\

\end{longtable}

\vspace{0.2cm}
{\small Lernziele, Begründung der Methodenwahl und offene Diskussionspunkte:
siehe \texttt{gewichtskraft-federkraft-\allowbreak lektionsplan.pdf}.}

\end{document}
